\chapter{Reprodução dos experimentos}
\label{cap:reproducao}

Este apêndice reúne as informações disponíveis para reconstruir os experimentos
descritos no livro: origem dos dados, ambiente computacional, hiperparâmetros e
protocolos de avaliação. Ele não equivale a um pacote de reprodução completo.

\section{Dados}

\subsection{O 3W}

O conjunto 3W é público e mantido pela Petrobras. Contém instâncias reais, simuladas e
desenhadas de eventos indesejáveis em poços de petróleo, rotuladas em dez classes
(\cref{tab:notb-classes}).

\begin{table}[htbp]
\centering
\caption{Procedência das instâncias do 3W e uso neste livro.}
\label{tab:rep-instancias}
\small
\begin{tabular}{lp{5.0cm}p{4.6cm}}
\toprule
\textbf{Tipo} & \textbf{Origem} & \textbf{Uso} \\
\midrule
\emph{WELL} & Poços reais &
Sempre incluído. Base de todos os resultados \\
\addlinespace
\emph{SIMULATED} & Simulador de escoamento &
Incluído apenas quando a classe tem menos de trinta instâncias reais \\
\addlinespace
\emph{DRAWN} & Traçado manualmente por especialista &
Sempre excluído \\
\bottomrule
\end{tabular}
\end{table}

\begin{destaque}[Por que excluir as instâncias desenhadas]
Instâncias \emph{DRAWN} são construções idealizadas do especialista. Elas representam o
que o evento \emph{deveria} parecer, não o que os sensores registraram.

Incluí-las eleva artificialmente as métricas, porque o modelo aprende a reconhecer a
idealização em vez do fenômeno. A exclusão é uma decisão de honestidade metodológica, e
custa desempenho reportado.
\end{destaque}

\subsection{Dados operacionais}

Os resultados do \cref{cap:dual} foram obtidos sobre dados de poços reais de um operador,
não públicos. Os três poços caracterizados naquele capítulo --- dois produtores e um
injetor com injeção alternada de água e gás, em lâminas d'água entre \num{1946} e
\num{2125} metros --- não podem ser disponibilizados.

As demais informações do capítulo --- arquitetura, hiperparâmetros, limiares e
protocolo --- são suficientes para reproduzir o método sobre outro conjunto.

\section{Ambiente computacional}

\begin{table}[htbp]
\centering
\caption{Ambiente computacional utilizado.}
\label{tab:rep-ambiente}
\small
\begin{tabular}{ll}
\toprule
\textbf{Componente} & \textbf{Especificação} \\
\midrule
Processador & Intel Xeon, 2 vCPU \\
Memória & \SI{13}{\giga\byte} \\
Aceleração & NVIDIA T4 \\
Linguagem & Python 3.7 (sistema em operação) e 3.10 (experimentos) \\
Aprendizado profundo & Keras sobre TensorFlow \\
Aprendizado clássico & scikit-learn \\
Orquestração em operação & Node-RED \\
Camada de linguagem & Qwen3.5-397B-A17B via NVIDIA NIM \\
\bottomrule
\end{tabular}
\end{table}

\begin{destaque}[O hardware é modesto de propósito]
Duas vCPU, treze gigabytes e uma T4 não constituem uma configuração de pesquisa
ambiciosa. É deliberado: um sistema destinado a operar em unidade flutuante precisa
caber em infraestrutura disponível a bordo ou em nuvem de custo controlado.

O ensemble completo ocupa \SI{21}{\mega\byte} e responde em \SI{405}{\milli\second}
(\cref{cap:closedworld}). Nenhum resultado do livro depende de escala computacional.
\end{destaque}

\section{Hiperparâmetros}

\subsection{Sistema dual (\cref{cap:dual})}

\begin{table}[htbp]
\centering
\caption{Configuração do autoencoder do sistema dual.}
\label{tab:rep-dual}
\small
\begin{tabular}{lp{9.4cm}}
\toprule
\textbf{Item} & \textbf{Valor} \\
\midrule
Entrada & $(500, 5)$ \\
Arquitetura & LSTM 16 $\to$ LSTM 4 $\to$ RepeatVector(500) $\to$ LSTM 16 $\to$
TimeDistributed Dense(5) \\
Compressão & $2500 \to 4$, fator $625$ \\
Perda & SMAPE \\
Otimizador & Adam \\
Épocas & 500 \\
Tamanho de lote & 100 \\
Ativação & ReLU \\
Normalização & Min-Max \\
Ruído de aumento & 0 a $0{,}5\,\%$ \\
Amostras de treinamento & $\approx 500$ janelas normais ($\approx \SI{5000}{\second}$) \\
Limiar & $\mu + 3\sigma$ sobre o SMAPE de treinamento \\
Níveis de risco & baixo $< 30\,\%$; médio $30$--$70\,\%$; alto $> 70\,\%$ \\
\bottomrule
\end{tabular}
\end{table}

\subsection{Ensemble binário (\cref{cap:binarios})}

\begin{table}[htbp]
\centering
\caption{Configuração dos classificadores binários um-contra-todos.}
\label{tab:rep-binarios}
\small
\begin{tabular}{lp{9.4cm}}
\toprule
\textbf{Item} & \textbf{Valor} \\
\midrule
Arquitetura & Entrada $(T,S)$ $\to$ LSTM 32--128 $\to$ Dropout 0,1--0,5 $\to$
Dense 32--64 ReLU $\to$ Dense 1 sigmoide \\
Busca de hiperparâmetros & Busca aleatória, 100 combinações por classe \\
Épocas & 25, com parada antecipada e paciência 5 \\
Tamanho de lote & 64 \\
Otimizador & Adam \\
Perda & Entropia cruzada binária \\
Partição & Temporal, $80/20$ \\
\bottomrule
\end{tabular}
\end{table}

\begin{destaque}[A partição temporal não é detalhe]
Partição aleatória sobre janelas deslizantes vaza informação: janelas consecutivas se
sobrepõem, e uma janela do conjunto de teste pode compartilhar amostras com uma de
treinamento.

A partição temporal --- primeiros $80\,\%$ para treinamento, últimos $20\,\%$ para teste
--- elimina o vazamento e produz números menores e honestos. Comparações com trabalhos
que usam partição aleatória são, por isso, desfavoráveis a este.
\end{destaque}

\subsection{Ciclo de mundo aberto (\cref{cap:mundoaberto})}

\begin{table}[htbp]
\centering
\caption{Configuração do ciclo de mundo aberto.}
\label{tab:rep-aberto}
\small
\begin{tabular}{lp{9.4cm}}
\toprule
\textbf{Item} & \textbf{Valor} \\
\midrule
Autoencoder & Entrada $(200,4)$ $\to$ LSTM 32 $\to$ LSTM 8 $\to$ Dense 16 (latente)
$\to$ LSTM 8 $\to$ LSTM 32 $\to$ Dense$(4 \times 200)$ \\
Compressão & $800 \to 16$, fator $50$ \\
Classificador latente & Dense 64 $\to$ Dropout 0,3 $\to$ Dense 32 $\to$ Dropout 0,2
$\to$ Dense 1 sigmoide \\
Limiar do voto & $\tau = 0{,}65$ \\
Estimador de $K$ & Floresta aleatória com oito atributos \\
Salvaguarda & Dois terços \\
\bottomrule
\end{tabular}
\end{table}

\subsection{Segmentação (\cref{cap:segmentacao})}

\begin{table}[htbp]
\centering
\caption{Configuração da U-Net unidimensional.}
\label{tab:rep-unet}
\small
\begin{tabular}{lp{9.4cm}}
\toprule
\textbf{Item} & \textbf{Valor} \\
\midrule
Entrada & $(L, 1)$, um modelo por sensor \\
Codificador 1 & Conv1D(32,5) + BN, $\times 2$; MaxPool(2); Dropout 0,2 \\
Codificador 2 & Conv1D(64,5) + BN, $\times 2$; MaxPool(2); Dropout 0,2 \\
Gargalo & Conv1D(128,5) + BN, $\times 2$ \\
Decodificador 1 & UpSampling(2); concatenação com 64; Conv1D(64,5) + BN, $\times 2$ \\
Decodificador 2 & UpSampling(2); concatenação com 32; Conv1D(32,5) + BN, $\times 2$ \\
Saída & Conv1D(1,1) sigmoide; limiar 0,5 \\
Janela & 200 passos $\times$ 4 variáveis; deslizamento de 100 passos, passo
\SI{5}{\second} \\
Pré-processamento & Imputação; padronização; remoção de ruído por ondaletas \\
Validação & Cruzada estratificada, cinco dobras \\
Otimizador & Adam; 30 a 50 épocas \\
\bottomrule
\end{tabular}
\end{table}

\subsection{Representação e Mahalanobis (\cref{cap:mahalanobis})}

\begin{table}[htbp]
\centering
\caption{Configuração dos espaços de representação avaliados.}
\label{tab:rep-mahal}
\small
\begin{tabular}{lp{9.4cm}}
\toprule
\textbf{Espaço} & \textbf{Definição} \\
\midrule
Bruto & 800 dimensões (200 passos $\times$ 4 variáveis) \\
Latente do autoencoder & 128 dimensões; Entrada $\to$ LSTM 256 $\to$ LSTM 128 $\to$
Dense 128 \\
Latente do multiclasse & 64 dimensões; penúltima camada do classificador \\
Cabeça binária & Latente(128) $\to$ Dense 64 ReLU $\to$ Dense 32 ReLU $\to$
Dense 2 \ing{softmax} \\
Limiar de novidade & Percentil 95 das distâncias de referência \\
Limiar binário & 0,3 \\
\bottomrule
\end{tabular}
\end{table}

\subsection{Camada de linguagem (\cref{cap:llm,cap:agente})}

\begin{table}[htbp]
\centering
\caption{Configuração da camada de linguagem.}
\label{tab:rep-llm}
\small
\begin{tabular}{lp{9.4cm}}
\toprule
\textbf{Item} & \textbf{Valor} \\
\midrule
Modelo & Qwen3.5-397B-A17B (mistura de especialistas) \\
Serviço & NVIDIA NIM \\
Modo de raciocínio & Ativado \\
Temperatura & 0,2 \\
Top-$p$ & 0,7 \\
Máximo de \ing{tokens} & 4096 \\
Formato de saída & JSON validado \\
Candidatas por consulta & 3, ordenadas \\
Pré-processamento & Escore-$z$; Min-Max $[0,1]$; ondaleta db4 nível 3, limiar suave;
tratamento de ausentes \\
Perfis de classe & Remoção de discrepantes por IQR em duas passagens, $k = 3{,}0$ e
depois $k = 1{,}5$ \\
\bottomrule
\end{tabular}
\end{table}

\begin{destaque}[O que não foi feito]
Sem ajuste fino. Sem exemplos no prompt. Sem múltiplas passagens com agregação. Sem
acesso a ferramentas externas.

Essas técnicas alterariam custo e desempenho em direções que precisam ser medidas; os
resultados do \cref{cap:resultados-agente} descrevem apenas a configuração documentada.
\end{destaque}

\begin{destaque}[Limites de reprodução desta edição]
O repositório desta edição não contém os scripts de treinamento e avaliação, as tabelas
derivadas em formato legível por máquina, as sementes das execuções, os ambientes
congelados nem os registros completos das chamadas ao serviço de linguagem. O
identificador exato do endpoint NIM e sua revisão também não estão documentados nas
fontes disponíveis. Assim, os parâmetros abaixo permitem reconstrução aproximada, mas
não repetição bit a bit nem auditoria independente de todos os números. Esses itens são
lacunas de evidência, não valores a serem inferidos.
\end{destaque}

\section{Protocolos de avaliação}

\begin{description}[leftmargin=0pt,style=nextline,itemsep=0.4em]

\item[Classe escondida (\cref{cap:closedworld}).] Treina-se o sistema com $|\mathcal{K}|
- 1$ classes, omitindo uma. Avaliam-se as instâncias da classe omitida: acerto significa
rejeição. Repete-se para cada classe.

\item[Comparação de arquiteturas (\cref{cap:mundoaberto}).] Mesma partição, mesmo
protocolo, mesma métrica; varia-se apenas o gerador da representação.

\item[Varredura de janela (\cref{cap:closedworld}).] Repete-se o protocolo de classe
escondida para comprimentos de janela de 150 a 7500 passos.

\item[Validação (\cref{cap:resultados-agente}).] Apresentam-se ao agente decisões
corretas e incorretas de um classificador; mede-se precisão e revocação no julgamento.

\item[Novidade com agente (\cref{cap:resultados-agente}).] Apresentam-se os perfis de
todas as classes exceto a do segmento; acerto significa declarar incompatibilidade.

\item[Intervalos de confiança.] Wilson a $95\,\%$ (\cref{eq:for-wilson}) quando as
contagens necessárias estão disponíveis; tabelas herdadas sem $n$ não permitem
reconstruí-los.

\item[Desvio entre poços.] Sempre que a avaliação é por poço, reporta-se média e desvio
padrão entre poços, e não apenas a média agregada.

\end{description}

\begin{destaque}[O item mais fácil de esquecer]
Reportar o desvio padrão entre poços foi o que tornou visível o principal ganho do
\cref{cap:segmentacao}: a acurácia subiu de $0{,}500$ para $0{,}863$, mas o desvio caiu
de $0{,}187$ para $0{,}082$.

Um sistema com desvio de $0{,}187$ entre poços exige investigação por poço antes de uma
decisão de implantação, ainda que sua média seja aceitável.
\end{destaque}

\section{Lista de verificação para reprodução}

\begin{enumerate}[leftmargin=*,itemsep=0.25em]
  \item Obter o 3W e verificar a versão utilizada.
  \item Excluir instâncias \emph{DRAWN}; incluir \emph{SIMULATED} apenas nas classes com
        menos de trinta instâncias reais.
  \item Usar partição temporal, nunca aleatória, sobre janelas deslizantes.
  \item Normalizar por poço, não globalmente.
  \item Fixar sementes e registrá-las.
  \item Reportar métricas por classe, não apenas a média global.
  \item Reportar desvio padrão entre poços.
  \item Usar intervalo de Wilson para todas as proporções.
  \item Registrar entradas e saídas da camada de linguagem, dado que o modelo servido por
        interface externa não é estritamente reprodutível.
  \item Documentar o prompt como código versionado.
  \item Exportar predições pareadas, contagens e tabelas derivadas em formato legível por
        máquina, junto ao script que gera cada tabela do livro.
  \item Congelar versões de dependências, identificador do modelo/endpoint e data de
        execução; registrar \ing{top-k}, semente e política de atualização quando o
        serviço os expuser.
\end{enumerate}
